\chapter*{Learning from the work}
\addcontentsline{toc}{chapter}{Introduction: Learning from the work}

A design record describes a change before or during implementation. A retrospective gives the work another opportunity to explain itself after encountering code, tests, review and users. It can expose an incorrect assumption, a boundary that held, a check that was insufficient or a useful simplification that was difficult to see at the start.

This part preserves those accounts alongside historical planning and abandoned work. The original paths identify their former roles. A retrospective, a graveyard entry and an old planning note answer different questions even when they concern the same feature.

\section*{Compare intention with outcome}

\subsection*{Compare the promise and the observation}

Read a retrospective with its linked design and integration evidence. Look for the difference between what was expected and what was observed. Which part of the original problem was solved? Which tradeoff was accepted? Did an unexpected failure change the design, or merely change the way it was checked?

\subsection*{Recover a reusable lesson}

An account is particularly useful when it names a concrete failure mode and the evidence that revealed it. Broad lessons become easier to reuse when a later reader can see the constraint, the attempted solution and the observed result together. A claim that a cycle was complete should retain the limits of its recorded checks.

\section*{Historical planning remains historical}

These documents were written under practices that changed over the project's life. Their cycle names, sponsor fields, lanes and ceremony record that context. They are preserved as evidence about the work, with current execution tracked outside the book.

An abandoned proposal may still contain a useful diagnosis. A successful change may still leave a debt or a later regression. Use this part to recover that knowledge and test its relevance to a new situation. Its purpose is to make the experience available for judgment, without treating an old plan as an automatically active work item.
